% draft_literature_review.tex
% Style: Academic Paper
% Target: 3--4 pages (~1,800--2,200 words of body text)

\section{Literature Review}
\label{sec:literature}

This section reviews three strands of literature that inform the present study: (1) investor
sentiment and asset pricing, (2) social media as a financial information channel, and (3) causal
identification in high-frequency panel settings. The paper sits at their intersection and
contributes a formal causal framework---Double Machine Learning---to a domain that has relied
primarily on reduced-form predictive regressions.

% ─────────────────────────────────────────────────────────────────────────────
\subsection{Investor Sentiment and Asset Pricing}
\label{sec:lit_sentiment}

The theoretical foundations for a sentiment channel in asset prices rest on limits to arbitrage.
When a subset of traders forms expectations using noisy signals---including affect, media tone, or
social contagion---prices can deviate from fundamental value if rational arbitrageurs face short
selling constraints, capital limits, or noise-trader risk \citep{delong1990, shleifer1997}.
\citet{baker2006} formalize this intuition, constructing a composite sentiment index from six
market-based proxies and showing that high-sentiment periods predict lower subsequent returns on
speculative stocks, consistent with initial overpricing followed by correction.

Empirically, the earliest text-based evidence comes from \citet{tetlock2007}, who applies the
Harvard-IV dictionary to the Wall Street Journal \textit{Abreast of the Market} column and finds
that media pessimism predicts next-day downward pressure on the Dow Jones Industrial Average, with
the effect fully reversing over the following week. \citet{tetlock2008} extend this to
firm-specific news and show that the negative word fraction in news stories predicts quarterly
earnings and stock returns even after controlling for prior earnings surprises. Together, these
papers establish that text sentiment carries incremental information beyond what standard financial
variables capture.

\citet{garcia2013} builds on Tetlock's framework by documenting time-varying effects: sentiment
extracted from the financial press is a stronger predictor of returns during recessions, when
information asymmetry and limits to arbitrage are likely more binding. \citet{loughran2011}
critique the use of general-purpose dictionaries in financial contexts, showing that the
Harvard-IV negative word list misclassifies domain-specific terms (e.g., ``liability,'' ``tax''),
and propose a finance-specific alternative that improves return prediction accuracy. This
measurement concern is directly relevant to the present study, which relies on Bloomberg's
proprietary dictionary whose construction is not disclosed.

% ─────────────────────────────────────────────────────────────────────────────
\subsection{Social Media as a Financial Information Channel}
\label{sec:lit_social_media}

A second strand of literature examines social media specifically as a venue for price-relevant
information transmission, distinct from traditional news outlets. The seminal contribution is
\citet{antweiler2004}, who analyze 1.5 million messages posted on Yahoo Finance and Raging Bull
message boards and find that message volume predicts next-day trading volume and volatility, with
a more modest link to returns. Their work establishes that user-generated financial discussion
contains informational content, but leaves open whether the effect is economically large or
causal.

\citet{bollen2011} shift the focus from message boards to Twitter. Using Google Profile of Mood
States (GPOMS) to extract six mood dimensions from public tweets, they report that the ``calm''
dimension Granger-causes changes in the DJIA with a lag of two to three days. Their approach is
aggregate---a single daily mood time series rather than firm-level variation---which limits
identification but demonstrates the existence of a predictive link between Twitter affect and
market movements.

\citet{chen2014} move to the firm level, examining posts on the Seeking Alpha investment platform.
They find that the tone of articles and, more strikingly, the tone of reader comments predicts
earnings surprises and subsequent abnormal returns, even after controlling for sell-side analyst
recommendations and traditional media tone. This suggests that the ``wisdom of crowds'' in social
media may reflect private information not yet incorporated in prices.

Two papers provide the most direct empirical antecedents for the present study.
\citet{gu2020} use firm-level daily Twitter sentiment from a commercial data provider to predict
open-to-open stock returns in a Fama-MacBeth framework over January 2015 to February 2017. Their
key results are: (i) a significant positive coefficient on one-day lagged Twitter sentiment
($+0.136^{***}$); (ii) a no-reversal test showing only lag~1 is significant when lags 1--5 are
entered jointly; and (iii) a profitable long-short strategy with an annualized Sharpe ratio of
3.17. These findings are interpreted as evidence that Twitter sentiment conveys firm-specific
information that is absorbed within one trading day. The present paper replicates the Gu and Kurov
specifications on a more recent sample (2024--2025) and finds substantially weaker effects,
suggesting possible temporal decay in the Twitter signal.

\citet{teti2019} examine Twitter data for 30 DJIA firms and find that the daily count of tweets
with negative polarity predicts stock returns. Their approach is pooled OLS with firm clustering,
and they introduce polarity-broken count variables (positive tweets, negative tweets, total
tweets) as an alternative to continuous sentiment scores. The present paper borrows their negative
tweet count construction (\texttt{twitter\_neg\_count}) and tests it alongside continuous
sentiment within the DoubleML framework.

More recent contributions expand the social media perimeter. \citet{cookson2023} use disagreement
across social media platforms (Twitter, StockTwits, Seeking Alpha) as a proxy for belief
dispersion and show that cross-platform disagreement predicts trading volume and return
volatility, suggesting that the informational content of social media is heterogeneous across
venues. \citet{bartov2018} find that aggregate Twitter opinion before quarterly earnings
announcements predicts announcement returns, suggesting that the Twitter crowd partially
anticipates earnings surprises. These findings collectively support the hypothesis that social
media contains price-relevant information, but they also highlight the difficulty of distinguishing
information from noise and causation from correlation.

% ─────────────────────────────────────────────────────────────────────────────
\subsection{Causal Identification in High-Frequency Panels}
\label{sec:lit_causal}

The studies reviewed above document predictive associations between sentiment measures and asset
returns. Whether these associations are causal is a distinct question that requires different
tools. The identification challenge is twofold: first, time-varying confounders---news events,
macroeconomic releases, institutional order flow---jointly drive both sentiment and prices;
second, reverse causality is plausible because large price movements generate social media
discussion, creating an endogenous feedback loop.

The standard approach to the first problem is panel fixed effects, which absorbs time-invariant
firm characteristics. This is the workhorse strategy in \citet{gu2020} and \citet{teti2019}, as
well as in the present paper's FE-OLS specifications. Fixed effects eliminate bias from stable
cross-sectional differences in sentiment coverage or trading clientele, but they do not address
time-varying confounders that are correlated with both sentiment and returns within a given firm.
Adding observable controls (news sentiment, momentum indicators, capital structure) partially
mitigates this concern, but only under the functional-form assumptions of OLS---specifically,
linearity in the confounder--outcome and confounder--treatment relationships.

The Double Machine Learning framework of \citet{chernozhukov2018} relaxes these functional-form
assumptions. The Partially Linear Regression (PLR) model assumes:
%
\begin{align}
    Y_{it} &= \theta_0 D_{it} + g_0(X_{it}) + U_{it}, \quad E[U_{it} \mid D_{it}, X_{it}] = 0 \label{eq:plr_outcome} \\
    D_{it} &= m_0(X_{it}) + V_{it}, \quad E[V_{it} \mid X_{it}] = 0 \label{eq:plr_treatment}
\end{align}
%
where $Y_{it}$ is the return, $D_{it}$ is the sentiment treatment, $X_{it}$ is the
high-dimensional confounder vector, and $\theta_0$ is the causal parameter of interest. The
nuisance functions $g_0(\cdot)$ and $m_0(\cdot)$ are estimated nonparametrically via machine
learning (here, XGBoost), and the Neyman-orthogonal score ensures that estimation error in the
nuisance functions does not contaminate inference on $\theta_0$ at $\sqrt{n}$ rates
\citep{chernozhukov2018}. Cross-fitting (sample splitting) prevents overfitting bias.

The DML framework has been applied in labor economics \citep{belloni2014}, program evaluation
\citep{knaus2021}, and demand estimation, but applications in empirical finance remain sparse.
\citet{fan2022} apply DML to predict bond returns with macroeconomic controls.
\citet{chetverikov2021} develop high-dimensional panel data methods in a related spirit. To my
knowledge, the present paper is the first to apply DoubleML specifically to intraday social media
sentiment and stock returns.

The second identification challenge---reverse causality---is more difficult. Instrumental
variables are a natural candidate, but credible instruments for social media sentiment are scarce:
potential instruments (e.g., platform outages, algorithmic feed changes) are rare, endogenous to
firm-level activity, or affect sentiment through channels that also affect prices directly. In the
absence of instruments, this paper adopts a placebo design: if sentiment is purely forward-looking
(information $\to$ prices), then today's sentiment should not predict yesterday's return. A large,
significant coefficient in this reverse regression is evidence of feedback---sentiment responding
to prior price movements---and constitutes a falsification test for the causal interpretation of
the forward estimates. This design follows the logic of Granger non-causality tests adapted for
panel settings, as in \citet{holtz-eakin1988}.

% ─────────────────────────────────────────────────────────────────────────────
\subsection{Positioning the Present Study}
\label{sec:lit_positioning}

The present paper sits at the intersection of these three literatures. It takes the research
question and variable constructions from the social media--finance literature (\citealt{gu2020};
\citealt{teti2019}), applies identification strategies from the causal inference literature
(\citealt{chernozhukov2018}), and contributes to the broader debate about whether investor
sentiment has real pricing effects (\citealt{baker2006}; \citealt{tetlock2007}). Three features
distinguish it from prior work:

\begin{enumerate}
    \item \textbf{Formal causal framework.} The DoubleML PLR estimator relaxes linearity
    assumptions that constrain prior FE and Fama-MacBeth estimates, providing a coefficient that
    is interpretable as a causal partial effect under the assumption that the confounder vector
    $X_{it}$ is sufficient for conditional exogeneity.

    \item \textbf{Reverse causality diagnostics.} Rather than assuming the direction of causation,
    the paper explicitly tests for it through a placebo regression and lag-structure analysis.
    The finding that the reverse coefficient ($+2.348$) exceeds the forward estimate ($-0.924$)
    in absolute value disciplines the causal interpretation throughout.

    \item \textbf{Updated sample period.} Prior Twitter-returns studies concentrate on 2011--2017,
    a period of rapid social media adoption. The present sample (2024--2025) captures a more
    mature informational environment in which algorithmic trading, bot activity, and information
    saturation may have attenuated the retail sentiment channel. The temporal comparison with
    \citet{gu2020} is informative about the stability of the Twitter effect over time.
\end{enumerate}


% =========================================================
% MISSING / EXPAND NOTES
% =========================================================
% [MISSING: .bib file. New keys required for this section: delong1990,
%   shleifer1997, baker2006, tetlock2008, garcia2013, loughran2011,
%   chen2014, bartov2018, belloni2014, knaus2021, fan2022,
%   chetverikov2021, holtz-eakin1988, cookson2023]
% [EXPAND: If the custom sentiment tool is further developed, add a
%   paragraph on entity-specific NLP approaches (FinBERT, domain-adapted
%   transformers) as a distinct methodological thread.]
% [EXPAND: Could add discussion of event-study approaches (e.g., Sprenger
%   et al. 2014 on S&P 500 tweets) if additional identification contrast
%   is desired.]
